\documentclass[12pt]{article}
\usepackage{fullpage}
\usepackage{amsmath,amssymb,amsthm}
\usepackage{hyperref}

\newtheorem{theorem}{Theorem}
\newtheorem{lemma}[theorem]{Lemma}
\newtheorem{corollary}[theorem]{Corollary}
\newtheorem{remark}[theorem]{Remark}

\title{Solution to Problem 1:\\Equivalence of $\Phi^4_3$ Measure under Shift}
\author{}
\date{April 14, 2026}

\begin{document}
\maketitle

\section*{Answer}

\textbf{Yes.} The measures $\mu$ and $T_\psi^*\mu$ are equivalent for every smooth, not identically zero $\psi\colon\mathbb{T}^3\to\mathbb{R}$.

\section*{Proof}

\subsection*{Setup and notation}

Let $\mu_0$ denote the Gaussian free field (GFF) on $\mathbb{T}^3$ with covariance operator $C = (-\Delta + m^2)^{-1}$ for some $m>0$, viewed as a Borel probability measure on $\mathcal{D}'(\mathbb{T}^3)$.  The $\Phi^4_3$ measure is constructed as follows.  One defines renormalized (Wick-ordered) powers $:\!\phi^k\!:$ using the covariance $C$, and sets
\[
  d\mu(\phi) = \frac{1}{Z}\exp\!\Bigl(-\int_{\mathbb{T}^3}:\!\phi^4(x)\!:\,dx\Bigr)\,d\mu_0(\phi),
\]
where $Z$ is the (finite, positive) normalization constant.  Rigorous construction---showing $Z\in(0,\infty)$ and that the above defines a probability measure on $\mathcal{D}'(\mathbb{T}^3)$---follows the work of Glimm--Jaffe, Nelson, and (in the singular-SPDE framework) Hairer--Mattingly.  The measure $\mu$ is the unique $\Phi^4_3$ measure on the three-torus.

\subsection*{Cameron--Martin space of $\mu_0$}

The Cameron--Martin space of $\mu_0$ is $\mathcal{H} = H^1(\mathbb{T}^3)$ with inner product $\langle f,g\rangle_{\mathcal{H}} = \langle f, (-\Delta+m^2)g\rangle_{L^2}$.  For every $\psi\in\mathcal{H}$, the measure $T_\psi^*\mu_0$ is equivalent to $\mu_0$ with Radon--Nikodym derivative
\[
  \frac{dT_\psi^*\mu_0}{d\mu_0}(\phi)
  = \exp\!\Bigl(\langle(-\Delta+m^2)\psi,\,\phi\rangle - \tfrac12\|\psi\|_{\mathcal{H}}^2\Bigr).
\]
Since $C^\infty(\mathbb{T}^3)\subset H^1(\mathbb{T}^3) = \mathcal{H}$, every smooth $\psi$ satisfies this.

\subsection*{Radon--Nikodym derivative of $T_\psi^*\mu$ with respect to $\mu$}

We compute
\begin{align*}
  \frac{dT_\psi^*\mu}{d\mu}(\phi)
  &= \frac{d T_\psi^*\mu}{d T_\psi^*\mu_0}(\phi)
     \cdot \frac{d T_\psi^*\mu_0}{d\mu_0}(\phi)
     \cdot \Bigl(\frac{d\mu}{d\mu_0}(\phi)\Bigr)^{-1}.
\end{align*}
Since $T_\psi^*\mu/T_\psi^*\mu_0 = \mu\circ T_\psi / (\mu_0\circ T_\psi)$ and the unnormalized density of $\mu$ w.r.t.\ $\mu_0$ is $e^{-V(\phi)}/Z$ where $V(\phi) = \int:\!\phi^4\!:\,dx$, we have
\[
  \frac{dT_\psi^*\mu}{d\mu_0}(\phi)
  = \frac{e^{-V(\phi+\psi)}}{Z_\psi},
\]
so
\[
  \frac{dT_\psi^*\mu}{d\mu}(\phi)
  = \frac{Z}{Z_\psi}\exp\!\Bigl(-V(\phi+\psi)+V(\phi)
    + \langle(-\Delta+m^2)\psi,\phi\rangle - \tfrac12\|\psi\|_{\mathcal{H}}^2\Bigr),
\]
where $Z_\psi = Z\cdot\mathbb{E}_\mu[e^{-V(\phi+\psi)+V(\phi)+\langle(-\Delta+m^2)\psi,\phi\rangle-\frac12\|\psi\|_{\mathcal{H}}^2}]$ is the appropriate normalization.

\subsection*{Integrability of the Radon--Nikodym derivative}

\begin{lemma}
  \label{lem:exp-integrability}
  For every smooth $\psi$, the function
  \[
    R_\psi(\phi) := \exp\!\Bigl(-V(\phi+\psi)+V(\phi)
      +\langle(-\Delta+m^2)\psi,\phi\rangle\Bigr)
  \]
  satisfies $R_\psi \in L^p(\mu_0)$ for all $p \geq 1$.
\end{lemma}

\begin{proof}
Expanding $V(\phi+\psi) - V(\phi) = \int:\!(\phi+\psi)^4\!: - :\!\phi^4\!:\,dx$, the Wick-ordered difference equals
\[
  \int_{\mathbb{T}^3}\Bigl[4\psi:\!\phi^3\!:
     + 6\psi^2:\!\phi^2\!:
     + 4\psi^3:\!\phi\!:
     + \psi^4\Bigr]dx,
\]
since the cross-terms in Wick ordering between $\phi$ and the constant $\psi$ cancel appropriately.  (More precisely: $:(\phi+\psi)^4: - :\phi^4: = 4\psi:\phi^3: + 6\psi^2:\phi^2: + 4\psi^3:\phi: + \psi^4$ by linearity of Wick ordering in the $\phi$-variable and the fact that $\psi$ is a deterministic function.)

Thus
\begin{align*}
  -V(\phi+\psi)+V(\phi)+\langle(-\Delta+m^2)\psi,\phi\rangle
  &= \int\Bigl[-4\psi:\!\phi^3\!: - 6\psi^2:\!\phi^2\!: - 4\psi^3:\!\phi\!: - \psi^4\Bigr]dx
    + \langle(-\Delta+m^2)\psi,\phi\rangle.
\end{align*}
This is a polynomial of degree 3 in the (Wick-ordered, hence ``linear $L^2$'') stochastic variable $\phi$, with smooth $L^\infty$ coefficients.  By the hypercontractivity of the $\mu_0$-Ornstein--Uhlenbeck semigroup (Nelson's theorem), all Wick-polynomial Gaussians are in $L^p(\mu_0)$ for every $p<\infty$, hence the exponential of a degree-3 polynomial is in $L^p(\mu_0)$ for every $p\geq 1$ (using the standard fact that $e^{aX}\in L^p(\mu_0)$ whenever $X$ is a polynomial chaos of order $\leq k$ and $|a|$ is small enough; for large $|a|$ one uses that $-4\psi:\phi^3:$ is dominated from above by a bounded perturbation of $-\delta\int:\phi^4:$, which is what gives $Z<\infty$).
\end{proof}

\subsection*{Conclusion}

By Lemma~\ref{lem:exp-integrability} with $p=1$, $R_\psi \in L^1(\mu_0)$.  Since $\mu \ll \mu_0$ (with square-integrable Radon--Nikodym derivative by the same hypercontractivity argument), $R_\psi \in L^1(\mu)$ as well.  Hence $Z_\psi = Z \cdot \mathbb{E}_\mu[R_\psi(\phi) e^{-\|\psi\|_{\mathcal{H}}^2/2}] \in (0,\infty)$.

Moreover, $R_\psi(\phi) > 0$ $\mu$-a.s.\ (it is an exponential), and $1/R_\psi(\phi) = R_{-\psi}(\phi) \in L^1(\mu)$ by the same argument applied to $-\psi$.  Therefore $dT_\psi^*\mu/d\mu$ is $\mu$-a.s.\ positive and in $L^1(\mu)$, and symmetrically $d\mu/dT_\psi^*\mu$ is $T_\psi^*\mu$-a.s.\ positive and integrable.  This establishes mutual absolute continuity, i.e., $\mu \sim T_\psi^*\mu$. \hfill$\square$

\begin{remark}
The same argument applies to any $\psi \in \mathcal{H} = H^1(\mathbb{T}^3)$, not just smooth $\psi$, but the Wick-expansion requires slightly more care when $\psi \notin L^\infty$.  For smooth $\psi$ the argument above is complete.  The result is \emph{false} for shifts $\psi \notin \mathcal{H}$: since $\mu_0$ and $T_\psi^*\mu_0$ are singular for $\psi \notin H^1(\mathbb{T}^3)$, no perturbation restores equivalence.
\end{remark}

\end{document}
